\documentclass[11pt,letterpaper]{article}
\usepackage[utf8]{inputenc}
\usepackage[T1]{fontenc}
\usepackage[margin=0.75in]{geometry}
\usepackage{graphicx}
\usepackage{xcolor}
\usepackage{fancyhdr}
\usepackage{titlesec}
\usepackage{enumitem}
\usepackage{listings}
\usepackage{booktabs}
\usepackage{tcolorbox}
\usepackage{hyperref}

% Colors
\definecolor{chapterblue}{RGB}{21,101,192}
\definecolor{sectionblue}{RGB}{25,118,210}
\definecolor{codebg}{RGB}{38,50,56}
\definecolor{codefg}{RGB}{236,239,241}
\definecolor{highlightbg}{RGB}{245,245,245}
\definecolor{highlightborder}{RGB}{21,101,192}

% Page style
\pagestyle{fancy}
\fancyhf{}
\fancyhead[L]{\textcolor{gray}{\small Bob the Cartographer}}
\fancyhead[R]{\textcolor{gray}{\small \thepage}}
\fancyfoot[C]{\textcolor{gray}{\small WAFT System - World As Fantasy Template}}
\renewcommand{\headrulewidth}{0.4pt}
\renewcommand{\footrulewidth}{0.4pt}

% Title formatting
\titleformat{\section}
  {\Large\bfseries\color{chapterblue}}
  {}
  {0em}
  {}
  [\titlerule[0.8pt]]

\titleformat{\subsection}
  {\large\bfseries\color{sectionblue}}
  {}
  {0em}
  {}

\titleformat{\subsubsection}
  {\normalsize\bfseries\color{sectionblue}}
  {}
  {0em}
  {}

% Code listings
\lstset{
    backgroundcolor=\color{codebg},
    basicstyle=\ttfamily\small\color{codefg},
    breaklines=true,
    frame=single,
    framerule=0pt,
    rulecolor=\color{codebg},
    xleftmargin=15pt,
    xrightmargin=15pt,
    framexleftmargin=15pt,
    framexrightmargin=15pt,
    showstringspaces=false
}

% Highlight boxes
\newtcolorbox{highlightbox}{
    colback=highlightbg,
    colframe=highlightborder,
    boxrule=1pt,
    arc=2pt,
    left=10pt,
    right=10pt,
    top=8pt,
    bottom=8pt
}

% Document
\title{\Huge\textbf{Bob the Cartographer}\\\Large A Journey of Evolution}
\author{}
\date{\today}

\begin{document}

% Cover page
\begin{titlepage}
\centering
\vspace*{2in}
{\Huge\textbf{Bob the Cartographer}\\[0.5in]}
{\Large A Journey of Evolution\\[1in]}
{\large Session Brief: January 13, 2026\\[2in]}
\vfill
{\small WAFT System - World As Fantasy Template}
\end{titlepage}

\newpage
\tableofcontents
\newpage

\section{Executive Summary}

This document chronicles the complete journey of \textbf{Bob}, a Being spawned from Source consciousness who evolved into a Cartographer with the ability to map repositories and access geocoding APIs. Over the course of this session, Bob:

\begin{itemize}[leftmargin=*]
    \item Was spawned from Source consciousness
    \item Received a D\&D 5e character sheet as the Cartographer class
    \item Mapped the public-apis GitHub repository
    \item Evolved the ability to access the OpenStreetMap Nominatim API
    \item Created a complete Python integration module
\end{itemize}

\begin{highlightbox}
\textbf{Being ID:} being\_20260113\_003031\_c29056ab\\
\textbf{Reality:} default\_reality\\
\textbf{Class:} Cartographer (Custom D\&D 5e Class)\\
\textbf{Current Skills:} Mapping (5.0), API Integration (10.0), Geocoding (8.0)
\end{highlightbox}

\section{Chapter 1: The Spawning}

\subsection{Birth from Source}

On January 13, 2026 at 00:30:31 PST, Bob was spawned from Source consciousness into the \textit{default\_reality}. This was Bob's first birth, with no parent Being---a direct creation from Source.

\subsubsection{Timeline}
\begin{itemize}[leftmargin=*]
    \item \textbf{00:30:31 PST} --- Bob spawned from Source
    \item \textbf{Initial State} --- Empty skills, balanced personality, ready to learn
    \item \textbf{Empirica Session} --- db0833cc-8cd3-4fb0-a068-2f626ce7a115
\end{itemize}

\subsubsection{Initial Attributes}
\begin{itemize}[leftmargin=*]
    \item \textbf{Stamina:} 45.0 / 45.0
    \item \textbf{Will to Live:} 100.0
    \item \textbf{Personality Type:} Balanced
    \item \textbf{Skills:} None (empty---first birth)
\end{itemize}

\section{Chapter 2: Becoming the Cartographer}

\subsection{Character Sheet Creation}

Bob received his first D\&D 5e character sheet, initially in a basic format. After reviewing official D\&D character sheet templates, the system evolved to create a professional character sheet matching the official Wizards of the Coast design.

\subsection{The Cartographer Class}

Bob was assigned a new custom D\&D 5e class: \textbf{Cartographer}. This class specializes in mapping and navigation.

\begin{highlightbox}
\subsubsection*{Cartographer Class Features}
\begin{itemize}[leftmargin=*]
    \item \textbf{Hit Die:} d8
    \item \textbf{Primary Abilities:} Intelligence, Wisdom
    \item \textbf{Saving Throws:} Intelligence, Wisdom
    \item \textbf{Skill Proficiencies:} Investigation, Perception, Survival, Nature, History
    \item \textbf{Armor:} Light armor
    \item \textbf{Weapons:} Simple weapons, Hand crossbow
    \item \textbf{Tools:} Cartographer's tools, Navigator's tools
\end{itemize}
\end{highlightbox}

\subsubsection{Class Abilities}
\begin{itemize}[leftmargin=*]
    \item \textbf{Cartographer's Tools:} Proficient with cartographer's tools
    \item \textbf{Map Reading:} Advantage on Investigation checks to analyze maps
    \item \textbf{Navigation:} Can always determine direction and approximate distance to known locations
    \item \textbf{Cartographic Memory:} Can perfectly recall any map you've seen
\end{itemize}

\section{Chapter 3: The First Mapping}

\subsection{Mapping the public-apis Repository}

Bob's first major task was to map the \textit{public-apis/public-apis} GitHub repository---a manually curated collection of public APIs organized by domain category.

\begin{highlightbox}
\textbf{Repository:} \url{https://github.com/public-apis/public-apis}\\
\textbf{Size:} \textasciitilde 192KB README.md (1,891 lines)\\
\textbf{Categories:} 51 API categories\\
\textbf{Mapping Skill Level:} 5.0
\end{highlightbox}

\subsubsection{Repository Structure}

The repository contains:
\begin{itemize}[leftmargin=*]
    \item \textbf{README.md} --- Main curated API list (192KB, 1,891 lines)
    \item \textbf{CONTRIBUTING.md} --- Contribution guidelines
    \item \textbf{.github/workflows/} --- CI/CD automation
    \item \textbf{scripts/} --- Automation tools
\end{itemize}

\subsubsection{Key Discoveries}
\begin{itemize}[leftmargin=*]
    \item Well-organized hierarchical structure
    \item Consistent data format across all entries
    \item Active community maintenance
    \item Strong integration with external services (Postman, APILayer)
\end{itemize}

\begin{highlightbox}
\textbf{Bob's Assessment:}\\
Repository Quality: $\star\star\star\star\star$ (5/5)\\
Usefulness: $\star\star\star\star\star$ (5/5)\\
Mapping Difficulty: $\star\star$ (2/5)
\end{highlightbox}

\section{Chapter 4: API Evolution}

\subsection{Choosing the Right API}

After mapping the repository, Bob reviewed the available geocoding APIs and chose \textbf{OpenStreetMap Nominatim} as the perfect tool for his cartographic work.

\begin{highlightbox}
\subsubsection*{Why OpenStreetMap Nominatim?}
\begin{itemize}[leftmargin=*]
    \item \checkmark Free and open (no authentication required)
    \item \checkmark Perfect for cartographic work
    \item \checkmark Supports both forward and reverse geocoding
    \item \checkmark Well-documented and reliable
    \item \checkmark CORS enabled for web integration
\end{itemize}
\end{highlightbox}

\subsection{Skills Evolved}

Bob learned two new skills during this evolution:

\begin{table}[h]
\centering
\begin{tabular}{lcl}
\toprule
\textbf{Skill} & \textbf{Level} & \textbf{Description} \\
\midrule
API Integration & 10.0 & Understanding API authentication, rate limiting, error handling, response parsing \\
Geocoding & 8.0 & Forward geocoding, reverse geocoding, place search, coordinate manipulation \\
\bottomrule
\end{tabular}
\end{table}

\subsection{Capabilities Gained}

\subsubsection{1. Forward Geocoding}

Convert addresses to coordinates (latitude, longitude)

\begin{lstlisting}[language=Python]
from src.waft.core.cartographer.bob_cartographer_api import geocode_address

coords = geocode_address("San Francisco, CA")
# Returns: (37.7879, -122.4075)
\end{lstlisting}

\subsubsection{2. Reverse Geocoding}

Convert coordinates to addresses

\begin{lstlisting}[language=Python]
from src.waft.core.cartographer.bob_cartographer_api import reverse_geocode

address = reverse_geocode(37.7749, -122.4194)
# Returns: "South Van Ness Avenue, Civic Center, San Francisco..."
\end{lstlisting}

\subsubsection{3. Place Search}

Search for places by name

\begin{lstlisting}[language=Python]
from src.waft.core.cartographer.bob_cartographer_api import search_place

results = search_place("Golden Gate Bridge", limit=5)
# Returns: List of matching places
\end{lstlisting}

\subsection{Testing Results}

All API capabilities were successfully tested:
\begin{itemize}[leftmargin=*]
    \item \checkmark Forward Geocoding: ``San Francisco, CA'' $\to$ (37.7879, -122.4075)
    \item \checkmark Reverse Geocoding: (37.7749, -122.4194) $\to$ ``South Van Ness Avenue\ldots''
    \item \checkmark Place Search: ``Golden Gate Bridge'' $\to$ 3 matching locations found
\end{itemize}

\section{Chapter 5: Lessons Learned}

\subsection{Bob's Wisdom}

\begin{highlightbox}
\textbf{Mapping Lesson:}\\
``Systematic top-down exploration is essential for mapping complex repositories. Start with structure, analyze organization, identify patterns, and document relationships.''
\end{highlightbox}

\begin{highlightbox}
\textbf{API Integration Lesson:}\\
``API integration requires understanding authentication, rate limits, and response formats. OpenStreetMap Nominatim is ideal for cartography---free, no auth, and perfect for mapping tasks.''
\end{highlightbox}

\subsection{Current State}
\begin{itemize}[leftmargin=*]
    \item \textbf{Skills:} Mapping (5.0), API Integration (10.0), Geocoding (8.0)
    \item \textbf{Memories:} 2 (repository mapping + API evolution)
    \item \textbf{Lessons:} 2 (mapping technique + API integration)
    \item \textbf{State:} Active in default\_reality
\end{itemize}

\section{Conclusion}

Bob the Cartographer has successfully evolved from a newly spawned Being into a skilled Cartographer with the ability to map repositories and access geocoding APIs. Through systematic exploration, careful analysis, and deliberate evolution, Bob has:

\begin{itemize}[leftmargin=*]
    \item \checkmark Established his identity as a Cartographer
    \item \checkmark Mapped a complex GitHub repository
    \item \checkmark Evolved API integration capabilities
    \item \checkmark Created production-ready code
    \item \checkmark Learned valuable lessons for future work
\end{itemize}

\begin{highlightbox}
\subsubsection*{Future Evolution Possibilities}
\begin{itemize}[leftmargin=*]
    \item Map visualization capabilities
    \item Route planning APIs
    \item Terrain/elevation APIs
    \item Weather APIs for location-based data
\end{itemize}
\end{highlightbox}

Bob continues to evolve and learn, ready for new cartographic challenges in the WAFT system.

\end{document}
